% Chapter 4 — bilingual terminology for AI-agent workflows (Hebrew--English BiDi).

\begin{chapterintrobox}
AI-agent workflows are often designed in English, while local teams may document decisions,
approvals, and operational policy in Hebrew. This chapter defines shared terminology for
automation, supervision, and real work across both languages \cite{mediumworkflowagents}.
\end{chapterintrobox}

\section{Why Bilingual Terminology Matters}
\label{sec:bidi-why}

English dominates vendor documentation, model APIs, and agent-framework vocabulary. Yet
organizations that automate real work must explain credentials, review gates, and escalation
rules to Hebrew-speaking operators and managers. Without aligned terms, teams misread what
an \emph{Agent}, \emph{Task}, or \emph{Process} actually authorizes in production
\cite{arxiv2506workagents}.

Live workforce research asks which occupational tasks workers are willing to delegate to
agents versus keep under direct human control \cite{arxiv2506workagents}. Bilingual
terminology is therefore not cosmetic---it is how mixed-language teams agree on boundaries
before automation scales \cite{recomaijobroles}.

\section{Core English--Hebrew Terminology}
\label{sec:bidi-glossary}

Each definition below lists the English term, Hebrew equivalent, and Meaning in workflow.
Definition rows avoid tabular BiDi conflicts while keeping bilingual pairs readable.

\vspace{0.35em}
\begin{itemize}
  \item \textbf{Agent} --- \texthebrew{סוכן} --- Executes a bounded task using tools.
  \item \textbf{Task} --- \texthebrew{משימה} --- A unit of work delegated to an agent.
  \item \textbf{Crew} --- \texthebrew{צוות} --- A coordinated group of agents or roles.
  \item \textbf{Process} --- \texthebrew{תהליך} --- Ordered workflow from input to output.
  \item \textbf{API Key} --- \texthebrew{מפתח גישה} ---
    Credential used to access external services.
  \item \textbf{Human-in-the-loop} --- \texthebrew{בקרה אנושית} ---
    Human review before high-risk output.
\end{itemize}

\section{Controlled Mixed-Language Communication}
\label{sec:bidi-mixed}

\begin{insightbox}{Mixed-language example}
\textbf{Mixed-language example.} In an Israeli engineering team, a design review might use
English terms such as Agent, Task, API Key, and Human-in-the-loop, while the operational
policy is explained in Hebrew.
\end{insightbox}

\noindent
The English labels above are standard in agent platforms. The operational policy itself can
be stated in Hebrew without embedding those Latin tokens inside Hebrew sentences.

\vspace{0.35em}
\begin{hebrew}
\begin{flushright}
בדיון פנימי, הצוות מפרט את מדיניות הפעולה בעברית.
\end{flushright}
\end{hebrew}

Practitioner forums note that complex delegation still fails when boundaries are vague
\cite{redditcomplextasks}. Workflow commentary stresses agents as connectors between tools
and teams---not silent replacements for accountability \cite{mediumworkflowagents}.

\label{sec:bidi-hebrew-summary}

\begin{tcolorbox}[colback=boxBg,colframe=boxFrame,title={\texthebrew{סיכום בעברית}},sharp corners]
\begin{hebrew}
\begin{flushright}
סוכני בינה מלאכותית יכולים לבצע משימות מוגדרות כחלק מתהליך עבודה מתוזמר.\\[0.45em]
במערכת כזו, המשימה מפורקת לשלבים קטנים וברורים.\\[0.45em]
הסוכן משתמש בכלים כדי לאסוף מידע, ליצור טיוטה או לבצע פעולה.\\[0.45em]
כאשר המשימה רגישה, מסוכנת או לא ברורה, נדרשת בקרה אנושית.\\[0.45em]
המטרה אינה להחליף בני אדם לגמרי, אלא להעביר משימות מובנות לאוטומציה תוך שמירה על אחריות, פיקוח ואיכות.\\[0.45em]
לכן חשוב להגדיר מראש מתי הסוכן פועל לבד ומתי האדם חייב לאשר את התוצאה.
\end{flushright}
\end{hebrew}
\end{tcolorbox}

\section{Operational Guidance for Mixed-Language Teams}
\label{sec:bidi-practical}

When Hebrew-English teams deploy agents into real work, three practices keep automation safe:
\begin{itemize}
  \item Publish the terminology definitions in design reviews so every role shares the same
    workflow vocabulary.
  \item Store credentials and API access rules in a single authoritative runbook language,
    with Hebrew summaries where operators require them.
  \item Mark customer-facing and compliance-sensitive steps as human-in-the-loop regardless
    of language \cite{redditcomplextasks,recomaijobroles}.
\end{itemize}

\begin{warningbox}{Accountability before scale}
Automation can accelerate structured tasks in any language context, but approval language
must be explicit before agents act on consequential work \cite{arxiv2506workagents}.
\end{warningbox}
